\section{Verification, source corrections, and further research}
\label{sec:audit}

\subsection{What the checks establish}

The mathematical proofs above are independent of floating-point
agreement. The supplied programs exercise three different layers of
the argument.

\paragraph{Exact finite algebra.}
The cubic symbolic script treats $\gamma,L,\zeta(2),\zeta(3)$,
the needed zeta derivatives, and $\Omega_3,J_3$ as formal symbols.
It reconstructs the four quadratic coefficients of $H$ from the
four stated restrictions and reduces the residual in
\eqref{ray:main} to the zero rational function. Its curve checks
compare direct truncated quotient expansion with
\eqref{curve:differentials}; they retain all path coefficients.
The mixed-order script independently constructs the original
coefficient convolution, the shuffle/partial-fraction formula, and
the output after eliminating nonpositive trailing orders. Its 624
exact rational checks include coincident colors and repeated poles.
A separate local-coefficient script checks the displayed Mellin
principal parts and Gamma jets and the finite harmonic
summation-by-parts algebra in Section~\ref{bridge:sec}. In total,
the three exact suites record 760 passed checks: 16 cubic/ray
identities, 624 mixed-order checks, and 120 local-coefficient checks.

\paragraph{A certified scalar evaluation.}
The code \texttt{verify\_cubic\_harmonic\_bridge.py} combines
Arb ball arithmetic with \eqref{bridge:tail-bound}, using
$N=50$, $K=24$ and 100 decimal digits of working precision.
The stored enclosure is
\[
 Q\in 1.9662627343915343406643753384116314238258689764916868267443076
       \ \pm\ 4.86\cdot10^{-62}.
\]
The analytic tail contribution alone is below $1.26\cdot10^{-62}$.
The reported enclosure includes rounding of its displayed center.
The corresponding enclosures for $\eta(1)$ and $\Omega_3$ are
obtained using the proved exact identities. They are not independent
proofs of those identities.

\paragraph{Independent numerical diagnostics.}
Ordinary high-precision quadrature of the regularized digamma cube
and the harmonic Mellin integral agree through approximately 70
decimal places. A different evaluator computes $\T$ by integrating
Jonqui\`ere's series on $(0,1)$ and an incomplete-Gamma tail on
$(1,\infty)$. Symmetric differences with polynomial extrapolation
check the third derivatives on six additional rays, including
$(1,-2,3)$. Their largest recorded absolute discrepancy is below
$5.7\cdot10^{-37}$. An independent computation of $I_2$ gives the
same second cubic coordinate through more than 30 decimal places.
These latter calculations use non-interval arithmetic and truncated
series; their digits are diagnostics, not certified error estimates.
The split at one is strictly inside the Jonqui\`ere radius $2\pi$.

The package records exact program versions, numerical parameters,
machine-readable outputs, and build instructions. No finite test
suite certifies the analytic convergence arguments: those arguments
are supplied in the text. The exact symbolic checks certify the
specified finite algebra, rather than all possible extensions of a
theorem.

\subsection{A source-status ledger for integration}

The source snapshot is pinned by its full commit identifier, with
blob hashes in \path{provenance/repository_inventory.json}.
All 80 manuscript chapter files and all 11 incoming archives matched
their recorded Git blob hashes. This establishes the identity of the
source corpus; it is not a claim that every proof in that corpus has
received a new exhaustive review. The mathematical audit concentrated
on the harmonic normalization, the cubic completion, the local
Tornheim decomposition, and the rational/colored reduction machinery
used by the present results.

\begin{longtable}{@{}>{\raggedright\arraybackslash}p{0.25\textwidth}>{\raggedright\arraybackslash}p{0.46\textwidth}>{\raggedright\arraybackslash}p{0.23\textwidth}@{}}
\toprule
Existing target & Result in this report & Status to record\\
\midrule
\endhead
Coincident report, Q1: third diagonal derivative
& Theorem~\ref{thm:diag-harmonic} and the absolutely convergent
harmonic series \eqref{bridge:arithmetic}.
& Reduced to $\eta(1)$; finite ordinary-constant closure remains open.\\[5pt]
Coincident report, Q2: asymmetric cubic directions
& Theorem~\ref{thm:cubic-rays}, with two fixed coordinates and
explicit coefficients; $\kappa$ also has the convergent integral
\eqref{ray:I2}.
& Complete spanning formula proved; arithmetic minimality unproved.\\[5pt]
Coincident report, Q11: mixed integer inner signs
& Theorem~\ref{mx:main}, including arbitrary integer weights,
coincident colors, repeated poles and all outer-order derivatives.
& Finite closure proved, with a guaranteed depth bound.\\[5pt]
Strict/inclusive harmonic conventions
& \eqref{eq:strict-inclusive} is retained exactly, including
the sign and the $\zeta'(2,a)$ shift.
& Existing normalization confirmed.\\[5pt]
Gaussian $S_6$ and revised $S_8$ candidates
& No exact target-vector certificate is supplied here.
& Remain open; no theorem label should be promoted.\\
\bottomrule
\end{longtable}

\subsection{Corrections and qualifications worth preserving}

No new false theorem was found in the specific repository results
used here. Several distinctions are nevertheless necessary for a
correct integration.

\begin{enumerate}
\item \textbf{Keep the full resonant numerators.}
The cubic completion \eqref{diag:completed} subtracts
$3h(t)^2/t$ and adds $3(1-t)\zeta(2-t)/(2t)$.
Replacing their numerators by their constant terms changes the
answer. This is an existing warning in the coincident report,
rederived here to protect the new harmonic bridge from a finite
normalization error.
\item \textbf{Use the mixed pole rule only in its stated domain.}
The rational sector has a finite simple-pole description. Its
extension with positive uncolored inner orders can have higher
pole order and infinitely many nonpositive poles. The explicit
$Z_{1,1,0}(C;1)$ example proves this difference. The earlier
rational theorem is not wrong; extending its pole statement without
the logarithmic correction would be wrong.
\item \textbf{Separate coordinate rank from arithmetic independence.}
Two formal quadratic coordinates of $H$ survive the exact slice
constraints, with slope coefficients $abc$ and $c\Delta$.
The theorem determines that reconstruction problem. It does not
show that $\eta(1)$ or $\kappa$ lies outside a chosen field of
ordinary constants, or that they are independent of each other.
\item \textbf{Restrict before differentiating.}
An admissible Tornheim ray is holomorphic at zero, while the full
three-variable function is meromorphic there. The values $1/3$ on
the diagonal and $5/12$ on the $C$ axis are consistent. A curve
tangent to $A+C=0$ or $B+C=0$ is outside
Theorem~\ref{thm:curve}; its value is not assigned by cancellation
of a zero denominator in the displayed ray formula.
\end{enumerate}

The incoming report already records corrections to the inspected
Borwein--Dilcher author PDF \cite{IncomingCoincident,BorweinDilcher}.
They should retain that attribution. In the integer Jonqui\`ere
formula, the convention is $H_0=0$, and the coefficient of
$x^{p-1}\log x$ is
\[
 B_p=\frac{(-1)^p}{\Gamma(p)},\qquad p\ge1,
\]
without an additional factor $H_{p-1}$. This follows at once from
\eqref{mx:positiveexp}; $p=3$ gives $-1/2$.
For the ordinary convergent local power series, the splitting point
must satisfy $0<\theta<2\pi$. These are previously documented
corrections to the inspected author version, not newly discovered
errors or assertions about every published version.

\subsection{Proposed further research questions}

The following problems are specific next steps. They are questions,
not conjectured equalities unsupported by calculation.

\paragraph{1. Determine the arithmetic content of the harmonic bridge.}
Fix a constant algebra before attempting reduction, for example one
generated over $\mathbb Q$ by ordinary Stieltjes constants, ordinary
zeta derivatives at integers, logarithms and specified Barnes values.
Does $\eta(1)$ have a finite expression in that algebra? The bridge
gives three analytic realizations of the same coordinate: an ordered
double-zeta Laurent coefficient, a convergent logarithmic harmonic
series, and a normalized digamma-cube integral. Seek a functional
identity relating two realizations without reintroducing the same
coefficient under another name.

\paragraph{2. Relate the second cubic coordinate to harmonic Laurent data.}
The constant $\kappa$ is a fixed asymmetric double-zeta derivative
and also the convergent moment \eqref{ray:I2} plus classical terms.
Can an ordered Hurwitz functional equation carry that derivative
from the origin to a harmonic Stieltjes coefficient at $(1,1)$?
A successful identity should state the exact finite correction and
the direction of continuation. The question of a relation between
$\kappa$ and $\eta(1)$ is distinct from the proved two-coordinate
spanning formula.

\paragraph{3. Classify useful cancellations among cubic directions.}
The conic $\Delta=0$ removes $\kappa$ from a single ray. More
generally, a finite combination $\sum_j r_jW_{v_j}'''(0)$ has no
displayed nonclassical coordinate whenever
$\sum_jr_ja_jb_jc_j=0$ and
$\sum_jr_jc_j\Delta_j=0$. Classify short rational configurations
with additional permutation, distribution, or character symmetry.
This is a finite exact linear problem and can produce concrete new
single-zeta identities without an independence assumption.

\paragraph{4. Construct the complete quartic directional layer.}
At fourth order the cubic homogeneous part of $H$ enters. Determine
its exact dimension after symmetry and the axis/stuffle restrictions,
then choose independent anchor slices and derive the analogue of
\eqref{ray:main}. The first goal is a proved finite coordinate
formula with convergent integral representatives. A finite
ordinary-zeta evaluation should be a separate theorem if achieved.

\paragraph{5. Extend the cubic bridge to a common Hurwitz shift.}
For $a>0$, the coefficient $\eta(a)$ has exact shift and derivative
relations in the ordered report. Apply the partial-fraction method
to shifted poles and compare with the primitive
$\mathcal P_3(a+1)-\mathcal P_3(a)$. Determine the explicit
boundary terms and whether products of shifted digamma and
polygamma functions give a useful identity involving $\eta(a)$.
The unit-coordinate finite part at a singular endpoint and an
ordinary integral over a positive shifted interval must stay distinct.

\paragraph{6. Obtain a general power-to-harmonic formula.}
For each $k\ge4$, the principal parts of $\psi(z)^k$ are finite
generalized-harmonic polynomials. The two-point contour argument
still suggests a normally convergent subtraction. Carry out the
finite summation-by-parts reduction explicitly for $k=4$ and identify
the required harmonic Laurent coefficients or higher-depth data.
The target is a compact identity with every constant fixed, rather
than an unspecified Mittag--Leffler entire remainder.

\paragraph{7. Find depth-reducing projections in the mixed sector.}
At root-of-unity colors, the finite reduction outputs a finite list
of nested functions and their first-order derivatives. Construct
exact matrices for character projections and distribution identities
on that list. Their kernels can identify combinations whose depth
drops to one, where Hurwitz-zeta, Gamma and generalized-Stieltjes
formulas become available. Any claimed minimal depth must specify
the relation space being used.

\paragraph{8. Evaluate selected holomorphic Laurent coefficients.}
Theorem~\ref{mx:poles} gives all principal parts but deliberately
does not promise finite depth-one formulas for the regular terms.
Begin with $Z_{1,1,0}(C;1)$ at $C=1,0,-1$: derive explicit
convergent kernels for its constant and linear Laurent coefficients
and relate them to depth-three harmonic Stieltjes conventions.
This separates the already solved local pole algebra from new
global arithmetic.

\paragraph{9. Normalize ordinary primitives under reflection and multiplication.}
The primitive $\mathcal P_3$ fixes its integration constant at the
singular origin, while \eqref{mx:primitives} fixes a holomorphic
Euler primitive at zero. Determine exact reflection and multiplication
relations for these normalizations, including their polynomial and
logarithmic correction terms. At low orders, compare with Barnes
multiple-Gamma primitives only after matching the additive constants
and branches.

\paragraph{10. Resolve tangential spectral paths.}
The transverse formula requires $(a+c)(b+c)\ne0$. For a curve
whose first nonzero contact with one polar divisor occurs at order
$r>1$, determine when the numerator cancels enough terms to leave
a removable restriction, and which path jets fix the resulting
finite value. This is a divisibility and coefficient-extraction
problem, not an instruction to apply the transverse formula at a
zero denominator.

\paragraph{11. Return to the exact Gaussian targets.}
The frozen $S_6$ and revised $S_8$ vectors still need exact
certificates in proved relation spaces. Use the current reduction
and projection machinery to propose additional rows, proving the
regularization prescription for each row. A numerical integer
relation is a candidate; a successful proof must combine exact
rows to the exact target. An unsuccessful restricted search should
report its precise space and a separating functional, not an
unrestricted impossibility claim.

\paragraph{12. Formalize the finite proof interfaces.}
The partial-fraction extraction, shuffle multiplicities, finite
antidifferences, Gamma jets, and cubic coefficient elimination are
all terminating algebra. Separate them from analytic lemmas for
normal convergence, Mellin continuation and endpoint constant
terms. A formal implementation can then check the finite identities
while explicitly declaring the analytic hypotheses on which their
applications depend.

These questions preserve the emphasis on exact identities and
normalizations. Error bounds appear in this report only to certify
the numerical evaluation of an already proved identity.
